\documentclass[11pt,a4paper]{article}
\usepackage[T1]{fontenc}
\usepackage{lmodern}
\usepackage{amsmath,amssymb,amsthm,mathtools}
\usepackage[margin=28mm]{geometry}
\usepackage{microtype}
\usepackage[hidelinks]{hyperref}
\usepackage{booktabs}
\newtheorem{theorem}{Theorem}
\newtheorem{lemma}[theorem]{Lemma}
\theoremstyle{remark}
\newtheorem{remark}[theorem]{Remark}
\DeclareMathOperator{\Des}{Des}
\DeclareMathOperator{\inv}{inv}
\newcommand{\NN}{\mathbb Z_{\geq0}}
\newcommand{\Sym}{\mathfrak S}
% Gaussian binomials use explicit standard matrix notation.
\newcommand{\wt}{\mathrm w}
\title{Strong $q$-log-concavity of $q$-descent coefficients}
\author{Xiaojian Zeng\\[2pt]
  \small PBC School of Finance, Tsinghua University, 100190\\
  \small \texttt{xiaojianzeng@foxmail.com}}
\date{}

\begin{document}
\maketitle
\begin{abstract}
We prove the strong coefficientwise $q$-log-concavity conjectured by
Gaetz and Gao for the coefficients of their $q$-descent expansion.
The proof establishes a directed inequality for Gaussian-binomial
differences with two independently chosen cutoffs. After coefficient
reversal, these kernels enumerate subsets whose minimum satisfies a
prescribed upper bound. We adapt Krattenthaler's shifted-set injection by
protecting the minimum of the larger subset, and give an explicit inverse
on the image. A direct decomposition of the defining permutations then
proves all the desired inequalities, including the diagonal cases, and
shows that every coefficient polynomial in the stated range is nonzero.
\end{abstract}

\section{The coefficients and the result}

Write $[n]=\{1,\ldots,n\}$, and let $q$ be an indeterminate. For
$\pi\in\Sym_n$, put
\[
 \Des(\pi)=\{h\in[n-1]:\pi_h>\pi_{h+1}\},\qquad
 \inv(\pi)=\#\{(h,\ell):h<\ell,\ \pi_h>\pi_\ell\}.
\]
For a finite nonempty $S\subset\mathbb Z_{>0}$, let $m=\max S$.
Following Gaetz--Gao \cite[Theorem 5]{GaetzGao}, define, for $1\leq k\leq m$,
\begin{equation}\label{eq:a-def}
 a_k(S;q)=
 \sum_{\substack{\pi\in\Sym_{m+k},\ \Des(\pi)=S\\
 [m+k]\setminus[m]\ \subseteq\ \{\pi_1,\ldots,\pi_m\}}}
 q^{\inv(\pi)},
 \qquad a_0(S;q)=a_{m+1}(S;q)=0.
\end{equation}
These are the coefficients in their expansion
\[
 \sum_{\substack{\pi\in\Sym_n\\\Des(\pi)=S}}q^{\inv(\pi)}
   =\sum_{k=1}^{m}a_k(S;q)\left[\begin{smallmatrix}n-m\\k\end{smallmatrix}\right]_{q}\qquad(n\geq m).
\]
Here $\left[\begin{smallmatrix}n\\k\end{smallmatrix}\right]_{q}$ is the Gaussian binomial coefficient, with value zero
outside $0\leq k\leq n$. We write $f\geq_q g$ when $f-g\in\NN[q]$.
A sequence has no internal zero polynomials if the indices of its
nonzero terms form an interval.

\begin{theorem}\label{thm:main}
For every finite nonempty $S\subset\mathbb Z_{>0}$, with $m=\max S$,
all the polynomials $a_k(S;q)$, $1\leq k\leq m$, are nonzero, and
\begin{equation}\label{eq:main}
 a_i(S;q)a_j(S;q)\geq_q a_{i-1}(S;q)a_{j+1}(S;q)
 \qquad(1\leq i\leq j\leq m).
\end{equation}
\end{theorem}

This proves the coefficient conjecture in
\cite[Conjecture 7]{GaetzGao}. The definition in Section 1.2 of the
consulted arXiv version prints $i<j$, while the nearby discussion also
uses $i=j$. Theorem~\ref{thm:main} proves the literal strict-index
statement and, explicitly, the additional diagonal inequalities.

Bencs proved the ordinary coefficient log-concavity at $q=1$
\cite[arXiv version, Proposition 4.3 and Corollary 4.4]{Bencs}.
Gaetz--Gao established strong $q$-log-concavity for a different
coefficient sequence, denoted $b_k(S;q)$ in their Theorem 6.
Their following remark shows that a general transformation of arbitrary
strongly $q$-log-concave $b$ sequences does not suffice for the $a$
conjecture \cite{GaetzGao}.
Our proof instead uses a decomposition into kernels for which a directed
inequality holds for every ordered pair of cutoffs.

The underlying transfer comes from Krattenthaler's proof of Gaussian
$q$-log-concavity \cite[Section 2]{Krattenthaler}. We retain its shifted
sets and last-maximum rule, but protect the minimum of the larger subset
before computing the prefix counts. This additional property permits
the two cutoff restrictions needed below.
Related work of Liang--Sagan proves log-concavity of descent counts
indexed by permutation length \cite[Theorem 6.6]{LiangSagan}; the index
in Theorem~\ref{thm:main} is instead the expansion index $k$.

\section{A directed inequality for cutoff kernels}

For $m\geq1$ and $1\leq k,r\leq m$, set
\begin{equation}\label{eq:K}
 K_{k,r}=q^{k^2}\left(\left[\begin{smallmatrix}m\\k\end{smallmatrix}\right]_{q}-\left[\begin{smallmatrix}m-r\\k\end{smallmatrix}\right]_{q}\right),
 \qquad K_{0,r}=K_{m+1,r}=0.
\end{equation}
The parameter $m$ is suppressed in this notation.

\begin{theorem}\label{thm:kernel}
For every $m\geq1$, $1\leq i\leq j\leq m$, and independently
$1\leq r,s\leq m$, one has
\begin{equation}\label{eq:kernel}
 K_{i,r}K_{j,s}-K_{i-1,r}K_{j+1,s}\in\NN[q].
\end{equation}
In particular, no ordering assumption on $r$ and $s$ is required.
\end{theorem}

For a $k$-subset $U$ of $[m]$, write
\[
 \wt_k(U)=\sum_{u\in U}u-\frac{k(k+1)}2,
 \qquad 0\leq\wt_k(U)\leq k(m-k).
\]
Let
\[
 \mathcal P_k(r)=\{U\subseteq[m]:|U|=k,\ \min U\leq r\},
 \qquad H_{k,r}(q)=\sum_{U\in\mathcal P_k(r)}q^{\wt_k(U)}
\]
for $1\leq k\leq m$. Define
$\mathcal P_0(r)=\mathcal P_{m+1}(r)=\varnothing$ and
$H_{0,r}=H_{m+1,r}=0$; no minimum of the empty subset is used.
The subset definition
\[
 \left[\begin{smallmatrix}n\\k\end{smallmatrix}\right]_{q}=\sum_{\substack{U\subseteq[n]\\|U|=k}}
 q^{\sum_{u\in U}u-k(k+1)/2}
\]
gives
\begin{equation}\label{eq:H}
 H_{k,r}(q)=\left[\begin{smallmatrix}m\\k\end{smallmatrix}\right]_{q}-q^{kr}\left[\begin{smallmatrix}m-r\\k\end{smallmatrix}\right]_{q}
 \qquad(1\leq k\leq m),
\end{equation}
since the excluded subsets are translations by $r$ of $k$-subsets of
$[m-r]$. Reflecting each element $u$ to $n+1-u$ proves Gaussian
reciprocity. Substituting it in \eqref{eq:H} yields
\begin{equation}\label{eq:reverse}
 q^{mk}H_{k,r}(q^{-1})=K_{k,r}(q).
\end{equation}
The equality also holds at the two zero boundaries. There are no negative
powers on its left. Moreover, $U=[k]$ gives the constant term $1$ in
$H_{k,r}$, so every interior kernel is nonzero.

\begin{lemma}[Transfer preserving a minimum]\label{lem:transfer}
Fix $1\leq i\leq j<m$ and put $\delta=j-i+1$. There is an injection
from pairs $(A,B)$ of subsets of $[m]$ of sizes $i-1,j+1$ to pairs
$(C,D)$ of sizes $i,j$ such that
\begin{align}
 \sum_{c\in C}c+\sum_{d\in D}d
   &=\sum_{a\in A}a+\sum_{b\in B}b-\delta,\label{eq:sum}\\
 \min D&=\min B,\qquad
 \min C\leq\min A\quad\text{if }A\ne\varnothing.\label{eq:min}
\end{align}
The two subsets in a pair need not be disjoint.
\end{lemma}

\begin{proof}
Protect $b=\min B$ by setting $B_0=B\setminus\{b\}$, and let
$A_\delta=\{a+\delta:a\in A\}$. For integer
$\delta\leq t\leq m+\delta$, define
\[
 f(t)=\#(A_\delta\cap[t])-\#(B_0\cap[t]).
\]
Every element of $A_\delta$ exceeds $\delta$, and $1\notin B_0$.
Consequently
\begin{equation}\label{eq:endpoint}
 f(\delta)\geq1-\delta>-\delta.
\end{equation}
For $t\geq m$, every element of $B_0$ has been counted, so
$f(t)\leq(i-1)-j=-\delta$.
Let $M$ be the maximum of $f$ on the indicated integer interval, and
let $L$ be the \emph{last} position at which this maximum is attained.
Then \eqref{eq:endpoint} implies $L<m$. Put $e=L+1$.
The increments of $f$ belong to $\{-1,0,1\}$; since $f(e)<f(L)$,
\[
 \delta<e\leq m,\qquad e\in B_0\setminus A_\delta,
 \qquad f(e)=M-1.
\]
Define
\begin{equation}\label{eq:map}
 C=A\cup\{e-\delta\},\qquad D=B\setminus\{e\}.
\end{equation}
The added element is in $[m]\setminus A$, and the removed element is
not $b$. Thus the sizes, \eqref{eq:sum}, and \eqref{eq:min} hold.

It remains to recover the input from an output in the image.
Recover $b$ as $\min D$, and on the same integer interval define
\[
 g(t)=\#\bigl((C+\delta)\cap[t]\bigr)
       -\#\bigl((D\setminus\{b\})\cap[t]\bigr).
\]
Equation~\eqref{eq:map} implies
\[
 g(t)=f(t)+2\,\mathbf1_{\{t\geq e\}}.
\]
For $t<e$, one has $g(t)\leq M$. For $t\geq e$, the definition of
the last maximum gives $f(t)\leq M-1$, hence $g(t)\leq M+1$;
equality holds at $t=e$. Therefore $e$ is the \emph{first} position
of the maximum of $g$. We recover
$A=C\setminus\{e-\delta\}$ and $B=D\cup\{e\}$.
This left inverse proves injectivity.
\end{proof}

\begin{proof}[Proof of Theorem~\ref{thm:kernel}]
Suppose first that $2\leq i\leq j<m$. Apply Lemma~\ref{lem:transfer}
to $A\in\mathcal P_{i-1}(r)$ and
$B\in\mathcal P_{j+1}(s)$. The minimum properties give
$C\in\mathcal P_i(r)$ and $D\in\mathcal P_j(s)$ independently
of the ordering of $r,s$. The normalization constants satisfy
\[
 \frac{(i-1)i}2+\frac{(j+1)(j+2)}2
 -\frac{i(i+1)}2-\frac{j(j+1)}2=\delta.
\]
Together with \eqref{eq:sum}, this shows that
\[
 \wt_i(C)+\wt_j(D)=\wt_{i-1}(A)+\wt_{j+1}(B).
\]
The injection is therefore weight preserving, and its image complement
enumerates the coefficients of
\begin{equation}\label{eq:H-ineq}
 H_{i,r}(q)H_{j,s}(q)-H_{i-1,r}(q)H_{j+1,s}(q)\in\NN[q].
\end{equation}
If $i=1$ or $j=m$, the subtracted product is zero by convention,
so \eqref{eq:H-ineq} holds in those cases as well, including $m=1$.
Apply $q^{m(i+j)}$ times substitution $q\mapsto q^{-1}$ to
\eqref{eq:H-ineq}. All exponents in both products are at most
$m(i+j)$, so this operation only reverses a finite coefficient list.
Equation~\eqref{eq:reverse} gives \eqref{eq:kernel}.
\end{proof}

\begin{remark}
The protection step has a concrete role. With $m=4$, $i=j=2$,
$A=\{2\}$ and $B=\{2,3,4\}$, the unmodified classical transfer
\cite[Section 2]{Krattenthaler} gives
$C=\{1,2\}$, $D=\{3,4\}$, losing the second minimum cutoff at $2$.
The map in Lemma~\ref{lem:transfer} gives $C=D=\{2,3\}$.
This is a distinction between maps, not a negative instance of a
polynomial inequality.
\end{remark}

\section{Recovering the descent coefficients}

\begin{proof}[Proof of Theorem~\ref{thm:main}]
Let $T=S\setminus\{m\}$. For $1\leq r\leq m$, define
\[
 x_r(q)=\sum_{\substack{\sigma\in\Sym_m\\
                   \Des(\sigma)=T,\ \sigma_m=r}}q^{\inv(\sigma)}.
\]
We prove directly from \eqref{eq:a-def} that
\begin{equation}\label{eq:reduction}
 a_k(S;q)=\sum_{r=1}^m x_r(q)K_{k,r}(q)
 \qquad(0\leq k\leq m+1).
\end{equation}
For $1\leq k\leq m$, a permutation counted in \eqref{eq:a-def} has
an increasing suffix $U=\{u_1<\cdots<u_k\}\subseteq[m]$.
Its prefix consists of $[m+k]\setminus U$; standardizing this prefix
gives a unique $\sigma\in\Sym_m$ with $\Des(\sigma)=T$.
If $\sigma_m=r$, then the descent at position $m$ is equivalent to
$r\geq u_1$, because exactly $u_1-1$ prefix values are smaller than
$u_1$. Conversely, any such $(\sigma,U)$ reconstructs one permutation
in \eqref{eq:a-def}: place the complement of $U$ in the relative order
$\sigma$, then append $U$ increasingly. Thus this is a bijection.

The prefix contributes $\inv(\sigma)$ inversions and the suffix
contributes none. For $u_t$, the number of larger prefix values is
$(m+k-u_t)-(k-t)=m-u_t+t$. Hence the cross inversions total
\[
 mk-\sum_{u\in U}u+\frac{k(k+1)}2=mk-\wt_k(U).
\]
Summing the weights and applying \eqref{eq:reverse} proves
\eqref{eq:reduction} for interior $k$. At the two boundaries it follows
from the zero conventions.

For $1\leq i\leq j\leq m$, expanding \eqref{eq:reduction} now gives
\begin{align*}
 a_i(S;q)a_j(S;q)&-a_{i-1}(S;q)a_{j+1}(S;q)\\
 &=\sum_{r,s=1}^m x_r(q)x_s(q)
       \bigl(K_{i,r}K_{j,s}-K_{i-1,r}K_{j+1,s}\bigr)\geq_q0
\end{align*}
by Theorem~\ref{thm:kernel} and $x_r\in\NN[q]$.

Finally, every $T\subseteq[m-1]$ is a descent set. To see this,
partition the positions into blocks joined by the edges in $T$,
assign consecutive value intervals to the blocks in increasing order,
and reverse the entries within each block. The resulting permutation
has descents exactly at $T$. Thus some $x_r$ is nonzero. For that $r$,
$U=[k]$ belongs to $\mathcal P_k(r)$ for every $1\leq k\leq m$.
Equation~\eqref{eq:reduction} therefore makes every interior $a_k$
nonzero, which also establishes the no-internal-zero condition.
\end{proof}

\section{Exact checks and scope of comparison}

The proof is uniform in all its parameters. The accompanying Python
verifier provides finite checks using integer arithmetic. It rebuilds
the subset families and defining permutations, and checks the transfer,
its inverse, and coefficient signs. This use of reconstructed integer
predicates and separate implementations has a methodological parallel
in the finite certificates of Zeng et al.
\cite[Sections 3--4]{ZengEtAl}: their verifier reconstructs the finite
graphs and crosschecks exact searches. Here the finite calculations
corroborate the displayed uniform proof.

The completed checks cover the following ranges:
\begin{center}
\begin{tabular}{@{}p{0.43\textwidth}p{0.49\textwidth}@{}}
\toprule
Check & Exhaustive range and count\\
\midrule
Transfer and inverse & All $1\leq i\leq j<m\leq11$;
  $1{,}444{,}125$ input pairs\\[3pt]
Directed kernel inequalities & All $1\leq i\leq j\leq m\leq16$,
  $1\leq r,s\leq m$; $112{,}676$ strict and $18{,}496$ diagonal
  cases for each of $H$ and $K$\\[3pt]
Original coefficient identities & Every nonempty $S$ with
  $\max S\leq5$ and every $1\leq k\leq\max S$; $129$ identities
  from $4{,}084{,}856$ enumerated permutations\\
\bottomrule
\end{tabular}
\end{center}
The kernel stage also checks $1{,}496$ pairs of the two subset/kernel
identities. Ordinary coefficient convolution crosschecks all $2{,}730$
directed $H$ inequalities through $m=7$.

For the larger checks, coefficient lists are packed into integers by
standard Kronecker substitution \cite[Section 3.1]{Harvey}.
Each factor counts at most $2^m$ subsets, so each product coefficient
is at most $4^m$. With base $B=2^{2m+4}$, adding a guard $B/2$ in
every digit before subtracting the two products leaves every digit
strictly between $0$ and $B$. There is no carry or borrow across digits;
the guard bit in a digit survives exactly when the corresponding
coefficient difference is nonnegative. This is a coefficient test,
not just a sign test at a positive value of $q$.
The source package includes the verifier, its results, and reproduction
instructions. These finite ranges do not by themselves establish an
inequality beyond the enumerated cases.

The closest inspected proof method is the shifted-set construction
in \cite{Krattenthaler}; the specific additional ingredient here is
preservation of the larger subset's minimum. Liang--Sagan's
$q$-order-ideal lemma \cite[Lemma 8.10]{LiangSagan} supplies another
general framework, but a specialization to the present cutoff kernels
has not been established here. The comparison with the full results of
Liu--Xiao on transformations preserving log-concavity \cite{LiuXiao}
remains incomplete because their full text was not accessible during
the source check of 13 September 2026. No first-proof or exhaustive
priority claim is made.

\paragraph{Computational and AI assistance.}
OpenAI Codex assisted with proof development, literature searches,
programming, separate exact checks, and manuscript preparation.
No formal proof assistant was used. The mathematical justification is
the explicit injection and permutation bijection given above; the
computations supply reproducible finite corroboration.

\bibliographystyle{plain}
\bibliography{references}
\end{document}
